\documentclass[11pt]{article}
\usepackage[a4paper,margin=2.2cm]{geometry}
\usepackage{newtxtext,newtxmath}
\usepackage[hidelinks]{hyperref}
\pagestyle{empty}
\setlength{\parskip}{6pt}
\begin{document}

\noindent
To the Editors,\\
\emph{International Journal of Heat and Mass Transfer}

\vspace{8pt}
\noindent
Dear Editors,

\noindent
Please consider our manuscript, ``Field accuracy is not heat accuracy:
wall-flux-constrained surrogates for steady two-dimensional non-Boussinesq convection'', for
publication in the
\emph{International Journal of Heat and Mass Transfer} as a full-length article.

Surrogates for thermal convection are usually trained on paired Boussinesq and compressible
solutions. Public paired data sets fix the relative temperature contrast, so a network trained
on them memorises one correction field and cannot rescale it at another contrast. The
manuscript asks whether the metric used to select these surrogates is the right one. On 200
held-out classical-cavity samples, a model with a 0.09\% median temperature-field error misses the
both-wall integral Nusselt number by 2.0\% on average (1.4--2.9\% across seeds); in 43\% of the
samples on average the field error is below 1\% and the $Nu$ error above 2\%. On a 22-sample
random-wall replication at $\varepsilon\approx1.2$, the corresponding field and $Nu$ errors are
0.47\% and 4.04\%, with 66\% on average in the same region. A wall-flux penalty in the training
loss reduces the both-wall error under the training six-point stencil to 0.01\% on the classical
cavity and 0.58\% in the replication, and the local wall-flux error to 0.08\% and 1.14\%; the
median field error moves from 0.09\% to 0.10\% and from 0.47\% to 0.49\%.
These six-point values are matched training/evaluation diagnostics against a non-converged
$64^2$ target, not absolute physical accuracy. The target correction carries a 0.43--0.60
percentage-point grid-refinement change, two orders of magnitude larger than 0.01\%, and a
stencil-free scaling model is worse than returning the Boussinesq field for $Nu$ at four, six,
eight and ten evaluation points. A complete two-seed $4\times4$ training/evaluation matrix
puts every row minimum on its training stencil. At ten points the six-point-trained models
give 1.00--2.21\%, against 1.11\% for returning $X$; the ten-point-trained models give
0.107--0.117\%. The gain is therefore estimator-specific and is not presented as a universal
physical improvement.

The underlying correction factorises as $D=\varepsilon P(X)+O(\varepsilon^2)$ with
$\varepsilon=\Delta T/T_0$. Across 627 paired solutions, the velocity correction has nominal
power-law exponents close to one, but resampling the 134 wall-mode families widens the
intervals. The quadratic curvature is positive for velocity and pressure, while its magnitude
and the temperature-curvature sign are unresolved. The prefactor is a functional of the flow.
Training a network on $P=D/\varepsilon$ and multiplying by $\varepsilon$ outside
makes a public single-contrast data set reusable within the tested boundary-condition family, reducing
the channel-averaged extrapolation error from 0.129 to 0.085. A plain model trained at the same
contrast and rescaled by $\varepsilon/\varepsilon_0$ performs within 0.2\% of the factorized
target, so applying $\varepsilon$ outside the learned correction, rather than the residual
parameterisation, is what makes the data reusable.

The three-dimensional evidence includes both a 33-pair exponent check and a separate 15-field
leave-one-contrast ROM evaluation. The 3D ROM improves the field floor by $1.26$--$2.40\times$,
but it uses one resolution and does not include a wall-flux target. Alternate geometry, fluid and
public-family tests are reported as consistency checks; none establishes general cross-family
or three-dimensional transfer.

The study is within the journal's scope: it quantifies the departure from the
Oberbeck--Boussinesq approximation, checks it on independently generated public data, and
provides a surrogate design that improves wall heat transfer under a matched estimator while
stating the remaining uncertainty. The main sweep is two-dimensional and steady; the 3D ROM
and exponent checks are reported separately, and the remaining scope is stated explicitly. The 256$^2$ Boussinesq solution is within 0.37\% of the
Le Qu\'er\'e benchmark at $Ra=10^6$. The work is original, is not under consideration
elsewhere, and all authors have approved the submission. The authors declare no competing
interests. The public data are cited, and our generated data, solver configuration and
reproduction scripts are available at the repository listed in the manuscript.

\vspace{4pt}
\noindent
Yours sincerely,\\
{[Corresponding author name]}\\
{[Affiliation]}\\
{[Email]}

\end{document}
